%\author{Author name}
\chapter{Lipsum Ipsum}
\label{ch01:chap01}

\chaptermark{Lipsum Ipsum}


\lipsum[5]

\section{Section Title}
\label{ch01:sec3.1}

\lipsum[5]

\begin{figure}
{\fcolorbox{grey}{grey}{\vbox to 10pc{\hbox to 33pc{\ }}}}
\caption{Fusce mauris. Vestibulum luctus nibh at lectus. Sed bibendum, nulla a faucibus semper, leo velit
ultricies tellu}
\end{figure}

\subsection{Sub Section Title}
\label{ch01:sec3.1.1}

\lipsum[5]

\begin{figure}
{\hspace*{-3pc}\fcolorbox{grey}{grey}{\vbox to 10pc{\hbox to 36pc{\ }}}}
{\hspace*{-3pc}\begin{minipage}{33pc}
\caption{Fusce mauris. Vestibulum luctus nibh at lectus. Sed bibendum, nulla a faucibus semper, leo velit
ultricies tellu}
\end{minipage}}
\end{figure}

The discharge coefficient is computed as:

\begin{example}{\label{ch01:exa3.1}}
\lipsum[5]

\solution \lipsum[5]
\lipsum[5]

\begin{nboxfig}
{\fcolorbox{grey}{grey}{\vbox to 10pc{\hbox to .5\textwidth{\ }}}}
{Fusce mauris. Vestibulum luctus nibh at lectus. Sed bibendum, nulla a faucibus semper, leo velit
ultricies tellu}
\end{nboxfig}
\end{example}

\lipsum[5]
\lipsum[5]

\subsubsection{C-heading}
\lipsum[5]

Consider a discretization of the scalar Lagrangian
\begin{equation}
\mathcal{L}(\phi)=\frac{1}{2}(\partial_\mu\phi)^2+V(\phi)(+J\cdot
\phi)\label{ch63:eqn63.1}
\end{equation}
and its path integral, on a (hyper)cubic lattice of size $a$, via
\begin{eqnarray}
&&x\rightarrow x_n;\;\;\; \phi(x)\rightarrow \phi_n=\phi(x_n);\;\;\;
\mathcal{D}\phi\rightarrow \prod_n d\phi_n;\nonumber\\
&&\partial_\mu\phi\rightarrow
\frac{1}{a}(\phi_{n+\mu}-\phi_n),\label{ch63:eqn63.2}
\end{eqnarray}
where $n+\mu$ is the nearest neighbor on the lattice (in direction
$\mu$) to site $n$.

\lipsum[5] the boundary cannot flow into the
watershed only due to gravity.

\paragraph{D-heading}
\lipsum[5]

\bigskip

\begin{equation}
\label{ch05:eqn5.9}
\tikzmarknode{C}{\mathrm{C}}_\mathrm{Q} = \frac{\tikzmarknode{Q}{\mathrm{Q}}_\mathrm{p}}{\mathrm{\tikzmarknode{n}{\mathrm{n}}\tikzmarknode{D}{D}}^3}
\end{equation}
%~\\
\begin{tikzpicture}[overlay,remember picture,
    color=voilet,
    >=Latex,node distance=.5cm,
    lbl/.style = {align=center, font=\footnotesize}  % <--- new
                        ]
\draw[<-] (C) -- ++ (-0.5,0)    node[lbl,left]        {discharge coefficient}; % <--- now the text is in two lines
\draw[<-] (Q) -- ++ (0.5,0.5)  node[lbl,above right]  {discharge capacity}; % node position is moved to left
\draw[<-] (n) -- ++ (-0.3,-0.5)  node[lbl,below left,left=-20pt,below=2pt] {rotation speed in revolution per unit time};% node position is moved to right
\draw[<-] (D) -- ++ (0.5,-0.1)  node[lbl,right]       {impeller diameter};
\end{tikzpicture}

\bigskip\smallskip

A watershed (Figure) is the\index{land area from which
water tends} to flow downward due to gravity along a drainage network
to a single common point, called the outlet.

\begin{enumerate}[(3)]
\item[(1)] Simple arithmetic average method. Simple arithmetic average method. Simple arithmetic average method

\item[(2)] Thiessen polygon area-weighted average method. Simple arithmetic average method.

\item[(3)] Thus, water inside the watershed cannot flow across its boundary due only to gravity; likewise, water outside the boundary cannot flow into the
watershed only due to gravity.
\end{enumerate}

\lipsum[5]



\lipsum[5]

\begin{table}
\processtable{Definition of hydrologic soil groups
(HSGs).\label{ch03:tab3.1}}
{\begin{tabular*}{\textwidth}{@{\extracolsep{\fill}}ll@{\raggedright}p{20pc}@{}}
\dtoprule
HSG & Runoff potential & \multicolumn{1}{l}{Definition} \\
\midrule
A & low & \multicolumn{1}{p{20pc}}{\raggedright High rate of water transmission; well to excessively drained sands or gravels; high infiltration rate, saturated 
hydraulic conductivity of the least transmissive layer $\textrm{K}_{\textrm{sat}} > 5.67\ \textrm{in\ hr}^{-1}$} \\[6pt]
B & moderately low &
\multicolumn{1}{p{20pc}}{\raggedright Moderate rate of water transmission; moderately well to well drained soils with moderately fine to moderately coarse 
textures; moderate infiltration rate, $\textrm{K}_{\textrm{sat}} = 1.42 \sim 5.67\ \textrm{in\ hr}^{ -1}$} \\[6pt]
C& moderately high&
\multicolumn{1}{p{20pc}}{\raggedright Slow rate of water transmission; soils with a layer that impedes downward movement of water or with moderately fine to fine 
texture; slow infiltration rate, $\textrm{K}_{\textrm{sat}} = 0.06 \sim 1.42\ \textrm{in\ hr}^{-1}$} \\[6pt]
D& high&
\multicolumn{1}{p{20pc}}{\raggedright Very slow rate of water transmission; clay soils, soils with a permanent high water table, soils with a claypan or clay 
layer at or near the surface, and shallow soils over nearly impervious material; very slow infiltration rate, $\textrm{K}_{\textrm{sat}} \le 0.06\ \textrm{in\ 
hr}^{ -1}$} \\
\dbottomrule
\end{tabular*}}
{\tabsource{The drainage network can\index{Zeta function!of an
endomorphism} consist of natural streams, manmade channels,
drainage ditches.}}
\end{table}

\lipsum[5]
\lipsum[5]
\begin{lemma}
\lipsum[5]
\end{lemma}

\lipsum[5]

\lipsum[5]


\lipsum[5]

\begin{theorem}{Sample Head}
\lipsum[5]
\end{theorem}

\lipsum[5]

\begin{definition}{Head text of Definition}
\lipsum[5]
\end{definition}

\lipsum[5]

\begin{chaptexbox}[Drainage network]
\lipsum[5]




\lipsum[5]
\end{chaptexbox}

\lipsum[5]

\begin{problemlist}{9.9}
\item \label{ch01:prob3.1} For the following USGS quadrangle
    map, identify and highlight the ridges and valleys or
    gullies.

\begin{figure}[!h]
\hspace*{18pt}\fcolorbox{grey}{grey}{\vbox to 10pc{\hbox to .5\textwidth{\ }}}
\end{figure}

\item \label{ch01:prob3.2} For the USGS quadrangle map of
    Problem, delineate the watersheds with their outlets at A
    and B, respectively. Also, determine the\index{geometric
    characteristics} of each of the two watersheds, including
    drainage area, longest flow path length, width, slope,
    slope length, and drainage\break density.

\item \label{ch01:prob3.3} The following topographic\index{map
    is superimposed} by 30-m grids to develop a digital
    elevation model (DEM).\index{balance examined} The elevations on the map have units
    of meters.
\begin{enumerate}[(3)]
\item[(1)] Develop the DEM

\item[(2)] Determine the flow directions using the D8 algorithm
manually

\item[(3)] Use subroutine \textit{FlowDir{\_}DEM( )} to check your
answer in (2)
\end{enumerate}

\item\label{ch01:prob3.15} For a site with loam soils, the
    results from a field infiltrometer test are shown in the
    following table.
\begin{enumerate}[(2)]
\item[(1)] Determine the\index{parameters} of the Horton and Philip models for
the site

\item[(2)] Estimate the actual\index{cumulative
    infiltration} and totalRiemann hypothesis!for
    curves\index{Riemann hypothesis!for curves} direct
    runoff resulting from a 1-hr storm with an average
    rainfall intensity of 45 mm hr$^{-1}$ using these two
    models

\bigskip

\begin{unnumtable} [!h]
\processtable{}
{\hspace*{-28pt}\begin{tabular}{@{}lrrrrrrrr}
\dtoprule Time (min)& 0& 2& 5& 10& 20& 30& 40&60 \\[3pt]
%%\midrule
Infiltration rate (mm hr$^{-1})$& 240& 140& 120& 102&
66&42& 27&27 \\
\dbottomrule
\end{tabular}}{}
\end{unnumtable}
\end{enumerate}


\item\label{ch01:prob3.16} ForWatershed and Water Balanc the
    site in Problem, if the
    initial soil moisture is 0.15, develop
    the\index{infiltration capacity} curve using the G-A model
    using the means of the soil water parameters listed in
    Table. Also, plot the developed curve versus the\index{Zeta function!Dedekind} measured
    values on a same graph to visually assess the performance
    of the G-A model.

\end{problemlist}


\begin{document}

%\author{Author name}
\chapter{Lipsum Ipsum}
\label{ch01:chap01}

\chaptermark{Lipsum Ipsum}


\lipsum[5]

\section{Section Title}
\label{ch01:sec3.1}

\lipsum[5]

\begin{figure}
{\fcolorbox{grey}{grey}{\vbox to 10pc{\hbox to 33pc{\ }}}}
\caption{Fusce mauris. Vestibulum luctus nibh at lectus. Sed bibendum, nulla a faucibus semper, leo velit
ultricies tellu}
\end{figure}

\subsection{Sub Section Title}
\label{ch01:sec3.1.1}

\lipsum[5]

\begin{figure}
{\hspace*{-3pc}\fcolorbox{grey}{grey}{\vbox to 10pc{\hbox to 36pc{\ }}}}
{\hspace*{-3pc}\begin{minipage}{33pc}
\caption{Fusce mauris. Vestibulum luctus nibh at lectus. Sed bibendum, nulla a faucibus semper, leo velit
ultricies tellu}
\end{minipage}}
\end{figure}

The discharge coefficient is computed as:

\begin{example}{\label{ch01:exa3.1}}
\lipsum[5]

\solution \lipsum[5]
\lipsum[5]

\begin{nboxfig}
{\fcolorbox{grey}{grey}{\vbox to 10pc{\hbox to .5\textwidth{\ }}}}
{Fusce mauris. Vestibulum luctus nibh at lectus. Sed bibendum, nulla a faucibus semper, leo velit
ultricies tellu}
\end{nboxfig}
\end{example}

\lipsum[5]
\lipsum[5]

\subsubsection{C-heading}
\lipsum[5]

Consider a discretization of the scalar Lagrangian
\begin{equation}
\mathcal{L}(\phi)=\frac{1}{2}(\partial_\mu\phi)^2+V(\phi)(+J\cdot
\phi)\label{ch63:eqn63.1}
\end{equation}
and its path integral, on a (hyper)cubic lattice of size $a$, via
\begin{eqnarray}
&&x\rightarrow x_n;\;\;\; \phi(x)\rightarrow \phi_n=\phi(x_n);\;\;\;
\mathcal{D}\phi\rightarrow \prod_n d\phi_n;\nonumber\\
&&\partial_\mu\phi\rightarrow
\frac{1}{a}(\phi_{n+\mu}-\phi_n),\label{ch63:eqn63.2}
\end{eqnarray}
where $n+\mu$ is the nearest neighbor on the lattice (in direction
$\mu$) to site $n$.

\lipsum[5] the boundary cannot flow into the
watershed only due to gravity.

\paragraph{D-heading}
\lipsum[5]

\bigskip

\begin{equation}
\label{ch05:eqn5.9}
\tikzmarknode{C}{\mathrm{C}}_\mathrm{Q} = \frac{\tikzmarknode{Q}{\mathrm{Q}}_\mathrm{p}}{\mathrm{\tikzmarknode{n}{\mathrm{n}}\tikzmarknode{D}{D}}^3}
\end{equation}
%~\\
\begin{tikzpicture}[overlay,remember picture,
    color=voilet,
    >=Latex,node distance=.5cm,
    lbl/.style = {align=center, font=\footnotesize}  % <--- new
                        ]
\draw[<-] (C) -- ++ (-0.5,0)    node[lbl,left]        {discharge coefficient}; % <--- now the text is in two lines
\draw[<-] (Q) -- ++ (0.5,0.5)  node[lbl,above right]  {discharge capacity}; % node position is moved to left
\draw[<-] (n) -- ++ (-0.3,-0.5)  node[lbl,below left,left=-20pt,below=2pt] {rotation speed in revolution per unit time};% node position is moved to right
\draw[<-] (D) -- ++ (0.5,-0.1)  node[lbl,right]       {impeller diameter};
\end{tikzpicture}

\bigskip\smallskip

A watershed (Figure) is the\index{land area from which
water tends} to flow downward due to gravity along a drainage network
to a single common point, called the outlet.

\begin{enumerate}[(3)]
\item[(1)] Simple arithmetic average method. Simple arithmetic average method. Simple arithmetic average method

\item[(2)] Thiessen polygon area-weighted average method. Simple arithmetic average method.

\item[(3)] Thus, water inside the watershed cannot flow across its boundary due only to gravity; likewise, water outside the boundary cannot flow into the
watershed only due to gravity.
\end{enumerate}

\lipsum[5]



\lipsum[5]

\begin{table}
\processtable{Definition of hydrologic soil groups
(HSGs).\label{ch03:tab3.1}}
{\begin{tabular*}{\textwidth}{@{\extracolsep{\fill}}ll@{\raggedright}p{20pc}@{}}
\dtoprule
HSG & Runoff potential & \multicolumn{1}{l}{Definition} \\
\midrule
A & low & \multicolumn{1}{p{20pc}}{\raggedright High rate of water transmission; well to excessively drained sands or gravels; high infiltration rate, saturated hydraulic conductivity of the least transmissive layer $\textrm{K}_{\textrm{sat}} > 5.67\ \textrm{in\ hr}^{-1}$} \\[6pt]
B & moderately low &
\multicolumn{1}{p{20pc}}{\raggedright Moderate rate of water transmission; moderately well to well drained soils with moderately fine to moderately coarse textures; moderate infiltration rate, $\textrm{K}_{\textrm{sat}} = 1.42 \sim 5.67\ \textrm{in\ hr}^{ -1}$} \\[6pt]
C& moderately high&
\multicolumn{1}{p{20pc}}{\raggedright Slow rate of water transmission; soils with a layer that impedes downward movement of water or with moderately fine to fine texture; slow infiltration rate, $\textrm{K}_{\textrm{sat}} = 0.06 \sim 1.42\ \textrm{in\ hr}^{-1}$} \\[6pt]
D& high&
\multicolumn{1}{p{20pc}}{\raggedright Very slow rate of water transmission; clay soils, soils with a permanent high water table, soils with a claypan or clay layer at or near the surface, and shallow soils over nearly impervious material; very slow infiltration rate, $\textrm{K}_{\textrm{sat}} \le 0.06\ \textrm{in\ hr}^{ -1}$} \\
\dbottomrule
\end{tabular*}}
{\tabsource{The drainage network can\index{Zeta function!of an
endomorphism} consist of natural streams, manmade channels,
drainage ditches.}}
\end{table}

\lipsum[5]
\lipsum[5]
\begin{lemma}
\lipsum[5]
\end{lemma}

\lipsum[5]

\lipsum[5]


\lipsum[5]

\begin{theorem}{Sample Head}
\lipsum[5]
\end{theorem}

\lipsum[5]

\begin{definition}{Head text of Definition}
\lipsum[5]
\end{definition}

\lipsum[5]

\begin{chaptexbox}[Drainage network]
\lipsum[5]




\lipsum[5]
\end{chaptexbox}

\lipsum[5]

\begin{problemlist}{9.9}
\item \label{ch01:prob3.1} For the following USGS quadrangle
    map, identify and highlight the ridges and valleys or
    gullies.

\begin{figure}[!h]
\hspace*{18pt}\fcolorbox{grey}{grey}{\vbox to 10pc{\hbox to .5\textwidth{\ }}}
\end{figure}

\item \label{ch01:prob3.2} For the USGS quadrangle map of
    Problem, delineate the watersheds with their outlets at A
    and B, respectively. Also, determine the\index{geometric
    characteristics} of each of the two watersheds, including
    drainage area, longest flow path length, width, slope,
    slope length, and drainage\break density.

\item \label{ch01:prob3.3} The following topographic\index{map
    is superimposed} by 30-m grids to develop a digital
    elevation model (DEM).\index{balance examined} The elevations on the map have units
    of meters.
\begin{enumerate}[(3)]
\item[(1)] Develop the DEM

\item[(2)] Determine the flow directions using the D8 algorithm
manually

\item[(3)] Use subroutine \textit{FlowDir{\_}DEM( )} to check your
answer in (2)
\end{enumerate}

\item\label{ch01:prob3.15} For a site with loam soils, the
    results from a field infiltrometer test are shown in the
    following table.
\begin{enumerate}[(2)]
\item[(1)] Determine the\index{parameters} of the Horton and Philip models for
the site

\item[(2)] Estimate the actual\index{cumulative
    infiltration} and totalRiemann hypothesis!for
    curves\index{Riemann hypothesis!for curves} direct
    runoff resulting from a 1-hr storm with an average
    rainfall intensity of 45 mm hr$^{-1}$ using these two
    models

\bigskip

\begin{unnumtable} [!h]
\processtable{}
{\hspace*{-28pt}\begin{tabular}{@{}lrrrrrrrr}
\dtoprule Time (min)& 0& 2& 5& 10& 20& 30& 40&60 \\[3pt]
%%\midrule
Infiltration rate (mm hr$^{-1})$& 240& 140& 120& 102&
66&42& 27&27 \\
\dbottomrule
\end{tabular}}{}
\end{unnumtable}
\end{enumerate}


\item\label{ch01:prob3.16} ForWatershed and Water Balanc the
    site in Problem, if the
    initial soil moisture is 0.15, develop
    the\index{infiltration capacity} curve using the G-A model
    using the means of the soil water parameters listed in
    Table. Also, plot the developed curve versus the\index{Zeta function!Dedekind} measured
    values on a same graph to visually assess the performance
    of the G-A model.

\end{problemlist}


\end{document}
